% Generated draft: requires mathematical and source review.
\documentclass[11pt]{article}
\usepackage[margin=25mm]{geometry}
\usepackage{fontspec}
\setmainfont{Latin Modern Roman}
\usepackage{amsmath,amssymb,mathtools,unicode-math}
\setmathfont{Latin Modern Math}
\usepackage{enumitem,needspace,xurl,hyperref}
\setlength{\parindent}{0pt}
\setlength{\parskip}{5pt}
\setlength{\emergencystretch}{3em}
\newcommand{\N}{\mathbb{N}}
\newcommand{\Z}{\mathbb{Z}}
\newcommand{\Q}{\mathbb{Q}}
\newcommand{\R}{\mathbb{R}}
\newcommand{\C}{\mathbb{C}}
\newcommand{\F}{\mathbb{F}}
\newcommand{\A}{\mathbb{A}}
\newcommand{\PP}{\mathbb{P}}
\newcommand{\E}{\mathbb{E}}
\newcommand{\eps}{\varepsilon}
\newcommand{\dd}{\,\mathrm{d}}
\begin{document}
\section*{UnsolvedMath ID: 30006336}
\textbf{Source problem code:} catalogue entry
\subsection{Definitions and mathematical statement}
\emph{Terminology.} A generic point of a topological space \(X\) is a point \(\eta\) such that \(\overline{\{\eta\}}=X\). This definition does not specify the geometric objects or cohomology theories intended in the source.

\emph{Requested extension, incompletely specified.} The registry asks whether generic-point vanishing for de Rham and étale cohomology extends to prismatic cohomology. The supplied report fragment states that Peter Scholze explained to the report's authors that their proof of Theorem 2 generalizes to prismatic cohomology, and explicitly says that they had not checked the details. The intended task is to formulate and establish the corresponding prismatic extension of that result.

The statement of Theorem 2 is absent from the supplied record. Consequently, the geometric hypotheses, base, coefficients, cohomological degrees, meaning of generic-point vanishing, and precise prismatic cohomology construction are unspecified. In particular, it is unclear whether vanishing concerns a cohomology group, a restriction map, or another construction. No fully quantified mathematical assertion can be recovered from this fragment without supplying those data. The reported expectation is not a completed theorem.
\subsection{Short English statement}
The intended question is whether the generic-point vanishing result described by the registry has a prismatic analogue and whether the proof referred to in the report extends to it. The authors explicitly left the details unchecked. The missing statement of Theorem 2 prevents a precise formulation of its hypotheses and conclusion.
\subsection{Sources}
\begin{itemize}
\item Local registry: \texttt{data/unsolvedmath/problems.json} and supplementary registry \texttt{catalog/problems.json}; permanent UnsolvedMath ID \(30006336\), exact source code \texttt{catalogue entry}, title \emph{Prismatic Extension of Generic-Point Vanishing}. The record assigns the year 2025 and names no formal proposer.
\item Supplied report citation: \emph{Non-Archimedean Geometry and Applications} (2025), Oberwolfach Reports, pp.\ 1517--1578, DOI \texttt{10.4171/owr/2025/29}. \url{https://doi.org/10.4171/owr/2025/29}. The supplied passage attributes the explanation about extending the proof to Peter Scholze; the report's authors are not identified in this record.
\item Supplied literature reference: \emph{Prisms and prismatic cohomology}. \url{https://arxiv.org/abs/1905.08229}. The supplied record does not identify a result in this reference that proves the requested extension.
\item Attribution and surrounding qualifications: the acknowledgements thank workshop participants, Marco D'Addezio, and Atsushi Shiho for discussions. They particularly thank Peter Scholze for discussions concerning the preceding section and his explanation of his programme over \(\Z\) relating de Rham and étale cohomology. They also thank Scholze and Alberto Vezzani for discussions of rigid cohomology in relation to Problem 5. Neither the preceding section nor Problem 5 is supplied; these acknowledgements do not specify hypotheses for the requested extension.
\item Supplied literature assessment, checked on \texttt{2026-08-22}: the record labels the problem open and describes the extension as an unchecked expectation rather than a completed theorem. This is supplied dated curation, not independently verified here. Its listed sources are the OWR report and \emph{Prisms and prismatic cohomology}, linked above; the record requests independent verification of time-sensitive status and cited claims.
\item Supplied statement-verification metadata: the registry labels the statement \texttt{corrected\_verified} and reports that the official report was checked, the question isolated from surrounding exposition, and notation repaired for self-contained publication. These are supplied provenance claims, not independent verification in this draft.
\end{itemize}
\end{document}
